\documentclass[10pt,a4paper]{amsart}
\usepackage[margin=19mm]{geometry}
\usepackage{amsmath,amssymb,mathtools}
\usepackage[scr=rsfs,cal=euler]{mathalfa}
\usepackage{xfrac}
\usepackage[unicode,hidelinks,bookmarks=false]{hyperref}
\newcommand{\ie}{i.e.\ }
\newcommand{\cf}{cf.\ }
\newcommand{\resp}{respectively\ }
\newcommand{\Subsection}{\subsection}
\providecommand{\nobreakdash}{\nobreak\hbox{-}}
\setlength{\emergencystretch}{2em}
\multlinegap0pt
\usepackage{bm}
\usepackage{algorithm}\usepackage{algpseudocode}
\newtheorem{lemma}{Lemma}
\renewcommand{\footnote}[1]{}
\title{SIMS: Scale-Invariant Merit-Function-Based Scalarization for Multi-Task Learning}
\author{Zebin Chen, Fei Xing, Yang Chen, Hua Liu, Andy HF Chow, Yuhua Qian, Yu Zhang}
\date{Source: arXiv:2609.12599v1 (CC BY 4.0)}
\begin{document}
\maketitle

\noindent\emph{Source and licence.} SIMS: Scale-Invariant Merit-Function-Based Scalarization for Multi-Task Learning. Version arXiv:2609.12599v1 (CC BY 4.0), distributed by arXiv under the Creative Commons Attribution 4.0 International licence, http://creativecommons.org/licenses/by/4.0/, as recorded in the arXiv licence metadata for this exact version and shown on the abstract page https://arxiv.org/abs/2609.12599v1. Full licence notice: \texttt{licenses/arxiv-2609.12599-CC-BY-4.0.txt}.\par\smallskip

\noindent\textit{Editorial scope.} Counted unit: the article's lemma lem:lse\_smoothness (source0.tex lines 2136--2152). The article's complete original proof of it is delivered with it (source0.tex lines 2153--2209, one \texttt{proof} environment of 57 lines). No mathematics is written, repaired or re-derived: the statement and the proof are the article's own lines, in the article's own notation, and no retained line is substituted or re-flowed.  No further source-local context is needed: every cross-reference of every retained line resolves inside the two delivered spans.  Nothing was omitted from any retained span, so the package carries no omission record.  No retained line contains a citation, so the wrapper carries no bibliography and no bibliography entry.  No retained line contains a figure environment, an image-inclusion command or drawing code, so no image is delivered and the package declares no asset dependency.  Editorial formatting only, in addition: an AMS \texttt{amsart} preamble that restates the article's own declarations and macros used by the retained lines, namely the environments \texttt{lemma} on separate counters, together with the standard packages those lines need.  The retained environments print in this document as Lemma 1 in order of appearance, whereas the article prints the same items with the numbering of its own full document.  The complete native source of the article as distributed in the arXiv e-print for this exact version is bundled unchanged as \texttt{sources/210133/source0.tex}.
\begin{lemma}[Smoothness of log-sum-exp smoothing]\label{lem:lse_smoothness}
Suppose that for all $i\in[m]$,
$\tilde{\mathcal{L}}_i$ is $\zeta_i$-smooth and $\ell_i$-Lipschitz continuous on $\mathbb{R}^d$.
Define
\begin{align}
    g_{\mu}(\theta)=\mu\ln\!\Big(\sum_{i=1}^{m}\exp\big(\tilde{\mathcal{L}}_i(\theta)/\mu\big)\Big),
    \quad \mu>0 \notag.
\end{align}
Then $g_{\mu}$ is continuously differentiable on $\mathbb{R}^d$, and its gradient is Lipschitz on $\mathbb{R}^d$ with
\begin{align}
    \|\nabla g_{\mu}(\theta)\!-\!\nabla g_{\mu}(\theta')\|
    \le\!
    \Big(\frac{\bar{\ell}^2}{\mu}+\bar{\zeta}\Big)
    \|\theta-\theta'\|,\, \forall~\theta,\theta'\in \!\mathbb{R}^d \notag.
\end{align} %\footnote{***cy:check here.}
Equivalently, $g_{\mu}$ is $\Big(\frac{\bar{\ell}^2}{\mu}+\bar{\zeta}\Big)$-smooth on $\mathbb{R}^d$, where $\bar{\ell}=\max_{i}\ell_i$ and $\bar{\zeta}=\max_{i}\zeta_i$.
\end{lemma}

\begin{proof}
Define $\boldsymbol{\tilde{\mathcal{L}}}(\theta):=(\tilde{\mathcal{L}}_1(\theta),\ldots,\tilde{\mathcal{L}}_m(\theta))^\top$ and
\[
h_\mu(z):=\mu\ln\Big(\sum_{i=1}^m e^{z_i/\mu}\Big),\qquad z\in\mathbb{R}^m,
\]
so that $g_\mu(\theta)=h_\mu(\boldsymbol{\tilde{\mathcal{L}}}(\theta))$. Let $p(z):=\nabla h_\mu(z)$ be the softmax vector, i.e.,
$p_i(z)=\exp(z_i/\mu)/\sum_{j}\exp(z_j/\mu)$, hence $\|p(z)\|_1=1$.
Moreover, $h_\mu$ is $(1/\mu)$-smooth:
\begin{equation}\label{eq:hmu_smooth_icml}
\|\nabla h_\mu(z)-\nabla h_\mu(z')\|\le \frac{1}{\mu}\|z-z'\|,\qquad \forall z,z'\in\mathbb{R}^m .
\end{equation}

By the chain rule, with $J_{\boldsymbol{\tilde{\mathcal{L}}}}(\theta)=(\nabla \tilde{\mathcal{L}}_1(\theta),\ldots,\nabla \tilde{\mathcal{L}}_m(\theta))$,
\[
\nabla g_\mu(\theta)=J_{\boldsymbol{\tilde{\mathcal{L}}}}(\theta)\,\nabla h_\mu(\boldsymbol{\tilde{\mathcal{L}}}(\theta))
=\sum_{i=1}^m p_i(\boldsymbol{\tilde{\mathcal{L}}}(\theta))\,\nabla \tilde{\mathcal{L}}_i(\theta).
\]
For any $\theta,\theta'\in \mathbb{R}^d$, we decompose
\begin{align*}
\|\nabla g_\mu(\theta)-\nabla g_\mu(\theta')\|
\le& \|J_{\boldsymbol{\tilde{\mathcal{L}}}}(\theta)\|\,\|\nabla h_\mu(\boldsymbol{\tilde{\mathcal{L}}}(\theta))-\nabla h_\mu(\boldsymbol{\tilde{\mathcal{L}}}(\theta'))\| \\ 
&+ \|(J_{\boldsymbol{\tilde{\mathcal{L}}}}(\theta)-J_{\boldsymbol{\tilde{\mathcal{L}}}}(\theta'))\,p(\boldsymbol{\tilde{\mathcal{L}}}(\theta'))\|.
\end{align*}
% \footnote{***cy:check here.}

For the first term, From Eq.~\eqref{eq:hmu_smooth_icml}, we have
\[
\|J_{\boldsymbol{\tilde{\mathcal{L}}}}(\theta)\|\,\|\nabla h_\mu(\boldsymbol{\tilde{\mathcal{L}}}(\theta))-\nabla h_\mu(\boldsymbol{\tilde{\mathcal{L}}}(\theta'))\|
\le \frac{\|J_{\boldsymbol{\tilde{\mathcal{L}}}}(\theta)\|}{\mu}\,\|\boldsymbol{\tilde{\mathcal{L}}}(\theta)-\boldsymbol{\tilde{\mathcal{L}}}(\theta')\|.
\]
Since each $\tilde{\mathcal{L}}_i$ is $\ell_i$-Lipschitz on $\mathbb{R}^d$, we have
$\|\nabla \tilde{\mathcal{L}}_i(\theta)\|\le \ell_i$ and $\|\boldsymbol{\tilde{\mathcal{L}}}(\theta)-\boldsymbol{\tilde{\mathcal{L}}}(\theta')\|\le (\max_i \ell_i)\|\theta-\theta'\|$.
Thus $\|J_{\boldsymbol{\tilde{\mathcal{L}}}}(\theta)\|\le \max_i \ell_i$, and the first term is bounded by
$\frac{\max_i \ell_i^2}{\mu}\|\theta-\theta'\|$.

For the second term, using $\|p(\boldsymbol{\tilde{\mathcal{L}}}(\theta'))\|_1=1$ and the $\zeta_i$-smoothness of $\tilde{\mathcal{L}}_i$,
\begin{align}
\|(J_{\boldsymbol{\tilde{\mathcal{L}}}}(\theta)\!-\!J_{\boldsymbol{\tilde{\mathcal{L}}}}(\theta'))\,p(\boldsymbol{\tilde{\mathcal{L}}}(\theta'))\|&=\!
\Big\|\sum_{i=1}^m p_i(\boldsymbol{\tilde{\mathcal{L}}}(\theta'))\big(\nabla \tilde{\mathcal{L}}_i(\theta)-\nabla \tilde{\mathcal{L}}_i(\theta')\big)\Big\| \notag\\
&\le \max_i \zeta_i \,\|\theta-\theta'\|\notag.
\end{align}
% \[
% \|(J_{\boldsymbol{\tilde{\mathcal{L}}}}(\theta)-J_{\boldsymbol{\tilde{\mathcal{L}}}}(\theta'))\,p(\boldsymbol{\tilde{\mathcal{L}}}(\theta'))\|
% =
% \Big\|\sum_{i=1}^m p_i(\boldsymbol{\tilde{\mathcal{L}}}(\theta'))\big(\nabla \tilde{\mathcal{L}}_i(\theta)-\nabla \tilde{\mathcal{L}}_i(\theta')\big)\Big\|
% \le \max_i \zeta_i \,\|\theta-\theta'\|.
% \]
% \footnote{***cy:check here.}
Combining the two bounds proves
\[
\|\nabla g_\mu(\theta)-\nabla g_\mu(\theta')\|
\le
\Big(\frac{\max_i \ell_i^2}{\mu}+\max_i \zeta_i\Big)\,\|\theta-\theta'\|,
\qquad \forall \theta,\theta'\in \mathbb{R}^d.
\]
Let $\bar{\ell}=\max_{i}\ell_i$ and $\bar{\zeta}=\max_{i}\zeta_i$. Then, $g_{\mu}$ is $\Big(\frac{\bar{\ell}^2}{\mu}+\bar{\zeta}\Big)$-smooth on $\mathbb{R}^d$.
\end{proof}

\end{document}
